% GENERATED FILE. Edit canonical lessons, then run npm run generate:books.
% canonical-source-hash: fnv1a64:b3882293ae2b6e94
% canonical-lessons: BN-C153-alochona-kora, BN-C153-proshongsa-kora, BN-C153-boka, BN-C153-dhonyobad-deoya, BN-C153-sohyo-kora

\chapter{To discuss, To praise, To scold, To thank, To put up with}
\label{ch:bn-to-discuss-to-praise-to-scold-to-thank-to-put-up-with}
\hlchaptermodality{153}

\begin{chapteropening}
\textbf{By the end of this chapter:} \emph{I can say to discuss, to praise, to scold, to thank and to put up with.}

Complete the last lesson of chapter 153: I can say to discuss, to praise, to scold, to thank and to put up with.
\end{chapteropening}

\section[\texorpdfstring{\bn{আলোচনা} \bn{করা} (ālochônā kôrā)}{আলোচনা করা (ālochônā kôrā)}]{\bn{আলোচনা} \bn{করা} (ālochônā kôrā) — to discuss}

\label{lesson:BN-C153-alochona-kora}



\begin{quote}
Before the new one: say the Bengali for to comb, then the Bengali for to bite.
\end{quote}

\subsection*{You'll want to know: \bn{আলোচনা} \bn{করা}}
\textbf{\bn{আলোচনা} \bn{করা}} — \emph{ālochônā kôrā} — \textquotedblleft{}to discuss\textquotedblright{}.

\begin{grammarlens}[title={What you've built}]
One more word for this chapter.
\end{grammarlens}

\subsection*{Your turn}
\begin{itemize}
\raggedright
  \item \emph{Say it:} \emph{ālochônā kôrā}
  \item \emph{Say it:} \emph{ālochônā kôrā}, once more
  \item \emph{Say it:} say \emph{ālochônā kôrā}
  \item \emph{Recall:} say the Bengali for to comb, then the Bengali for to bite, then say \emph{ālochônā kôrā} again
\end{itemize}

\begin{tcolorbox}[breakable,colback=teal!4,colframe=teal!35!black,title={Before you move on}]
What does \bn{আলোচনা} \bn{করা} mean? (\textquotedblleft{}To discuss\textquotedblright{}.) Say it once more.
\end{tcolorbox}
\section[\texorpdfstring{\bn{প্রশংসা} \bn{করা} (prôshôngsā kôrā)}{প্রশংসা করা (prôshôngsā kôrā)}]{\bn{প্রশংসা} \bn{করা} (prôshôngsā kôrā) — to praise}

\label{lesson:BN-C153-proshongsa-kora}



\begin{quote}
Before the new one: say the Bengali for to bite, then the Bengali for to discuss.
\end{quote}

\subsection*{You'll want to know: \bn{প্রশংসা} \bn{করা}}
\textbf{\bn{প্রশংসা} \bn{করা}} — \emph{prôshôngsā kôrā} — \textquotedblleft{}to praise\textquotedblright{}.

\begin{grammarlens}[title={What you've built}]
One more word for this chapter.
\end{grammarlens}

\subsection*{Your turn}
\begin{itemize}
\raggedright
  \item \emph{Say it:} \emph{prôshôngsā kôrā}
  \item \emph{Say it:} \emph{prôshôngsā kôrā}, once more
  \item \emph{Say it:} say \emph{prôshôngsā kôrā}
  \item \emph{Recall:} say the Bengali for to bite, then the Bengali for to discuss, then say \emph{prôshôngsā kôrā} again
\end{itemize}

\begin{tcolorbox}[breakable,colback=teal!4,colframe=teal!35!black,title={Before you move on}]
What does \bn{প্রশংসা} \bn{করা} mean? (\textquotedblleft{}To praise\textquotedblright{}.) Say it once more.
\end{tcolorbox}
\section[\texorpdfstring{\bn{বকা} (bôkā)}{বকা (bôkā)}]{\bn{বকা} (bôkā) — to scold}

\label{lesson:BN-C153-boka}



\begin{quote}
Before the new one: say the Bengali for to discuss, then the Bengali for to praise.
\end{quote}

\subsection*{You'll want to know: \bn{বকা}}
\textbf{\bn{বকা}} — \emph{bôkā} — \textquotedblleft{}to scold\textquotedblright{}.

\begin{grammarlens}[title={What you've built}]
One more word for this chapter.
\end{grammarlens}

\subsection*{Your turn}
\begin{itemize}
\raggedright
  \item \emph{Say it:} \emph{bôkā}
  \item \emph{Say it:} \emph{bôkā}, once more
  \item \emph{Say it:} say \emph{bôkā}
  \item \emph{Recall:} say the Bengali for to discuss, then the Bengali for to praise, then say \emph{bôkā} again
\end{itemize}

\begin{tcolorbox}[breakable,colback=teal!4,colframe=teal!35!black,title={Before you move on}]
What does \bn{বকা} mean? (\textquotedblleft{}To scold\textquotedblright{}.) Say it once more.
\end{tcolorbox}
\section[\texorpdfstring{\bn{ধন্যবাদ} \bn{দেওয়া} (dhônyôbād deoyā)}{ধন্যবাদ দেওয়া (dhônyôbād deoyā)}]{\bn{ধন্যবাদ} \bn{দেওয়া} (dhônyôbād deoyā) — to thank}

\label{lesson:BN-C153-dhonyobad-deoya}



\begin{quote}
Before the new one: say the Bengali for to praise, then the Bengali for to scold.
\end{quote}

\subsection*{You'll want to know: \bn{ধন্যবাদ} \bn{দেওয়া}}
\textbf{\bn{ধন্যবাদ} \bn{দেওয়া}} — \emph{dhônyôbād deoyā} — \textquotedblleft{}to thank\textquotedblright{}.

\begin{grammarlens}[title={What you've built}]
One more word for this chapter.
\end{grammarlens}

\subsection*{Your turn}
\begin{itemize}
\raggedright
  \item \emph{Say it:} \emph{dhônyôbād deoyā}
  \item \emph{Say it:} \emph{dhônyôbād deoyā}, once more
  \item \emph{Say it:} say \emph{dhônyôbād deoyā}
  \item \emph{Recall:} say the Bengali for to praise, then the Bengali for to scold, then say \emph{dhônyôbād deoyā} again
\end{itemize}

\begin{tcolorbox}[breakable,colback=teal!4,colframe=teal!35!black,title={Before you move on}]
What does \bn{ধন্যবাদ} \bn{দেওয়া} mean? (\textquotedblleft{}To thank\textquotedblright{}.) Say it once more.
\end{tcolorbox}
\section[\texorpdfstring{\bn{সহ্য} \bn{করা} (sôhyô kôrā)}{সহ্য করা (sôhyô kôrā)}]{\bn{সহ্য} \bn{করা} (sôhyô kôrā) — to put up with}

\label{lesson:BN-C153-sohyo-kora}



\begin{quote}
Before the new one: say the Bengali for to scold, then the Bengali for to thank.
\end{quote}

\subsection*{You'll want to know: \bn{সহ্য} \bn{করা}}
\textbf{\bn{সহ্য} \bn{করা}} — \emph{sôhyô kôrā} — \textquotedblleft{}to put up with\textquotedblright{}.

\begin{grammarlens}[title={What you've built}]
One more word for this chapter.
\end{grammarlens}

\subsection*{Your turn}
\begin{itemize}
\raggedright
  \item \emph{Say it:} \emph{sôhyô kôrā}
  \item \emph{Say it:} \emph{sôhyô kôrā}, once more
  \item \emph{Say it:} say \emph{sôhyô kôrā}
  \item \emph{Recall:} say the Bengali for to scold, then the Bengali for to thank, then say \emph{sôhyô kôrā} again
\end{itemize}

\begin{tcolorbox}[breakable,colback=teal!4,colframe=teal!35!black,title={Before you move on}]
What does \bn{সহ্য} \bn{করা} mean? (\textquotedblleft{}To put up with\textquotedblright{}.) Say it once more.
\end{tcolorbox}
